\documentclass[oneside,12pt]{scrartcl}
\usepackage[utf8]{inputenc}
\usepackage{graphicx}
\usepackage{eso-pic}
\usepackage{caption}
\usepackage{subcaption}
\usepackage{minted}
\usepackage{xcolor}
\usepackage{silence}
\usepackage[hidelinks]{hyperref}
\usepackage{arydshln}
\usepackage{mathtools}
\usepackage[toc]{glossaries}
\usepackage{parskip}
\usepackage{tikz}  
\usepackage[T1]{fontenc}




\newcommand*{\escape}[1]{\texttt{\textbackslash#1}}
\newcommand{\insererText}{\textcolor{red}{[INSERER TEXTE ICI]}}
\newcommand{\insererTextT}[1]{\insererText\ \textcolor{red}{#1}}
\newcommand{\insererImage}{\textcolor{red}{[INSERER IMAGE ICI]}}
\newcommand{\HRule}{\rule{\linewidth}{0.5mm}} % Defines a new command for the horizontal lines, change thickness here


\usepackage{xpatch,letltxmacro}
\LetLtxMacro{\cminted}{\minted}
\let\endcminted\endminted
\xpretocmd{\cminted}{\RecustomVerbatimEnvironment{Verbatim}{BVerbatim}{}}{}{}

\renewcommand{\thesection}{\Roman{section}} %redefini apparence nom sections (affichage chiffres romains)

\RedeclareSectionCommand[
  tocnumwidth=1cm
]{section}
\RedeclareSectionCommand[
  tocnumwidth=1.1cm
]{subsection}

\title{COMPTE-RENDU XXXX}
\subtitle{XXXX \\ M1 Informatique}
\author{AUTEUR}
\date{XX/XX/2022}

\AddToShipoutPicture*
{
\put(235,100) %(x=100, y=100)
 {\includegraphics[scale=0.3]{imgs/logo-univ-rouen.png}}
}

\begin{document}

\begin{titlepage}

\center % Center everything on the page

%----------------------------------------------------------------------------------------
%	TITLE SECTION
%----------------------------------------------------------------------------------------

%\HRule \\[0.4cm]
%{ \huge \bfseries Interaction between Artificial Intelligence and Computer Architecture}\\[0.4cm] % Title of your document
%\HRule \\[1.5cm]

\HRule \\[0.4cm]
{ \huge \bfseries Protocole 2D-Doc et Algorithmes}\\[0.4cm] % Title of your document
%{ \huge \bfseries Leveraging Artificial Intelligence Algorithms \\ to Design Hardware Prediction Mechanisms}\\[0.4cm] % Title of your document
\HRule \\[1.5cm]

\textsc{\huge \textbf{Documentation Technique}}\\[1.5cm]

\textsc{\Large \textbf{Université de Rouen Normandie \\[0.07cm] UFR des Sciences et Technique \\[1,5cm] Master 1 Informatique\\[0.5cm] Projet Datamatrix M1 SSI 2021-2022}}


%----------------------------------------------------------------------------------------
%	AUTHOR & ADVISOR(S) SECTION
%----------------------------------------------------------------------------------------

\vspace{3cm}

\begin{center} \large
\Large \emph{Auteur:}\\[0.1cm]
\large BELABDOUN Billal\\
\end{center}
\vspace{1cm}

%----------------------------------------------------------------------------------------
%	DATE SECTION
%----------------------------------------------------------------------------------------

{\large Mai 2022}\\[1cm]
%----------------------------------------------------------------------------------------

\vfill % Fill the rest of the page with whitespace

\end{titlepage}
\thispagestyle{empty} %pour retirer la 1ere page de la pagination

\newpage

\begin{center}
\textsc{\Large \textbf{Révision du document}}
\end{center}
\vspace{2cm}
\begin{center}
\begin{tabular}{|c|c|c|}
    \hline
  N° de version & Date de modification & Modifications apportées  \\
    \hline
    1 & 12/05/2022 & Création du document\\
    \hline
\end{tabular}
\end{center}

\newpage
\pagenumbering{arabic} %remet la pagination pour le sommaire
{
\renewcommand*\contentsname{Sommaire}
\tableofcontents
}

\clearpage


\printglossary[title=Glossaire,nonumberlist]



\section{Introduction}
De nos jours la sécurisation des serveurs est indispensable, car les cybercriminels les ciblent en priorité. En effet, c’est l’endroit idéal pour y trouver des données sensibles. 

Lors de l'exploitation d'un serveur Web, il est important de mettre en place des mesures de sécurité pour protéger votre site et vos utilisateurs. Vous pouvez protéger vos sites Web et vos applications avec, par exemple, des politiques de pare-feu et restreindre l'accès à certaines zones avec une authentification par mot de passe. Ceci est un excellent point de départ pour sécuriser votre système, mais n'est pas suffisant. Il est important de s’assurer que des mesures de sécurité robustes sont intégrées à votre serveur et vos applications afin de maintenir un niveau de sécurité élevé. \\
Avec la pandémie et le télétravail, les problèmes de piratage ont nettement augmenté. Pourtant, il existe un certain nombre de mesures pouvant être prises pour garantir un niveau élevé de sécurité. Voici une liste non exhaustive de mesures pouvant être prises facilement afin d'augmenter la sécurité d'un serveur : 
\begin{itemize}
    \item Établir  \hyperref[POLITIQUE]{une politique de mots de passe}
    \item Sécuriser votre serveur internet avec \hyperref[HTTPS]{HTTPS}
    \item Maintenir à jour votre serveur
    \item Préparer le serveur contre les attaques par force brute (Fail2ban)
    \item Changer le port SSH par défaut
\end{itemize}

\end{document}